\newpage
\section*{Resumen}
Determina el alcance del trabajo. Con la lectura de este texto el lector debe enterarse claramente de qué se trata el trabajo. Debe ser informativo del desarrollo del mismo. No debe incluir juicios de valor, interpretaciones o críticas.\par
El Resumen debe contar con MENOS de 400 palabras. 
\vspace{11pt}

\section*{Área Temática}
Computación e Informática, Digitales.
\vspace{11pt}

\section*{Asignaturas}
Informática, Métodos numéricos, Probabilidad y estadística, Electrónica digital III.
\vspace{11pt}

\section*{Palabras Claves}
Osciloscopio. Impulso tipo rayo. Alta tensión. 
\vspace{11pt}

\newpage
\section*{Abstract}
Es la redacción en idioma Inglés del Resumen. 
\vspace{11pt}

\section*{Thematic area}
Computación e Informática, Digitales.
\vspace{11pt}

\section*{Subjects}
Informática, Métodos numéricos, Probabilidad y estadística, Electrónica digital III.
\vspace{11pt}

\section*{Key Words}
Es la traducción al idioma Inglés de las Palabras Claves.
\vspace{11pt}

\newpage
\section*{Resumo}
Es la redacción en idioma Portugués del Resumen.
\vspace{11pt}

\section*{Área Temática}
Computación e Informática, Digitales.
\vspace{11pt}

\section*{Assuntos}
Informática, Métodos numéricos, Probabilidad y estadística, Electrónica digital III.
\vspace{11pt}

\section*{Palavras Chaves}
Es la traducción al idioma Portugués de las Palabras Claves.
\vspace{11pt}







		




 








